\Needspace{8\baselineskip}
\Article{282}
Finding the representations of a given binary form whose determinant is not
$=0$\HistoricalNote{*)}{For brevity, we pass over this case here; it requires a somewhat different method.}
by a given ternary form depends on the following observations:

I.\ From any proper representation of the binary form $(p,q,r)=\varphi$ of determinant
$D$ by the ternary form $f$ of determinant $\Delta$, we can derive
integers $B,B'$ such that
\[
BB\equiv\Delta p,\qquad BB'\equiv-\Delta q,\qquad
B'B'\equiv\Delta r\ (\mathrm{mod.}\ D)
\]
\textit{i.e.}, a value of the expression $\surd\Delta(p,\allowbreak -q,\allowbreak r)(\mathrm{mod.}\ D)$.
Let a proper representation of the form $\varphi$ by $f$ be given by
\[
x=\alpha t+\beta u,\qquad x'=\alpha't+\beta'u,\qquad x''=\alpha''t+\beta''u
\]
(where $x,x',x''$; $t,u$ denote the indeterminates of the forms $f,\varphi$).
Choose integers $\gamma,\gamma',\gamma''$ so that
\[
(\alpha'\beta''-\alpha''\beta')\gamma+
(\alpha''\beta-\alpha\beta'')\gamma'+
(\alpha\beta'-\alpha'\beta)\gamma''
\]
$=k$ is either $=+1$ or $=-1$, and let $f$ be transformed by the substitution
\[
\begin{array}{ccc}
\alpha,&\beta,&\gamma\\
\alpha',&\beta',&\gamma'\\
\alpha'',&\beta'',&\gamma''
\end{array}
\]
into the form
$\left(\begin{smallmatrix}a,a',a''\\b,b',b''\end{smallmatrix}\right)=g$,
whose adjoint is
$\left(\begin{smallmatrix}A,A',A''\\B,B',B''\end{smallmatrix}\right)=G$.
Then clearly $a=p$, $b''=q$, $a'=r$, $A''=D$, and $\Delta$ will be the
determinant of the form $g$, so
\[
BB=\Delta p+A'D,\qquad BB'=-\Delta q+B''D,\qquad B'B'=\Delta r+AD
\]
Thus, \textit{for example}, the form $19tt+6tu+41uu$ is represented by $xx+x'x'+x''x''$
on setting $x=3t+5u$, $x'=3t-4u$, $x''=t$. Then, setting $\gamma=-1$, $\gamma'=1$, $\gamma''=0$,
\SourceLeaf{329}
we obtain $B=-171$, $B'=27$, or the value $(-171,27)$ of the expression
$\surd-1(19,\allowbreak -3,\allowbreak 41)(\mathrm{mod.}\ 770)$.

It follows that, if $\Delta(p,-q,r)$ is not a quadratic residue of $D$,
$\varphi$ cannot be properly represented by any ternary form of determinant
$\Delta$. So, in the case where $\Delta,D$ are relatively prime, $\Delta$
must be a characteristic number of the form $\varphi$.

II.\ Since $\gamma,\gamma',\gamma''$ can be chosen in infinitely many different ways,
different values of $B,B'$ will result. Let us see how they are related.
Suppose that $\delta,\delta',\delta''$ too are such that
\[
(\alpha'\beta''-\alpha''\beta')\delta+
(\alpha''\beta-\alpha\beta'')\delta'+
(\alpha\beta'-\alpha'\beta)\delta''=\mathfrak{k}
\]
is either $=+1$ or $=-1$, and that the form $f$ is transformed
by the substitution
\[
\begin{array}{ccc}
\alpha,&\beta,&\delta\\
\alpha',&\beta',&\delta'\\
\alpha'',&\beta'',&\delta''
\end{array}
\]
into
$\left(\begin{smallmatrix}\mathfrak{a},\mathfrak{a}',\mathfrak{a}''\\\mathfrak{b},\mathfrak{b}',\mathfrak{b}''\end{smallmatrix}\right)=\mathfrak{g}$,
whose adjoint is
$\left(\begin{smallmatrix}\mathfrak{A},\mathfrak{A}',\mathfrak{A}''\\\mathfrak{B},\mathfrak{B}',\mathfrak{B}''\end{smallmatrix}\right)=\mathfrak{G}$.
Then $g,\mathfrak{g}$ will be equivalent, and so will $G$ and $\mathfrak{G}$.
Applying the rules given in Articles 269, 270\HistoricalNote{*)}{From the transformation of the form $f$ into $g$, derive the transformation of $g$ into $f$; from this and the transformation of $f$ into $\mathfrak{g}$, derive the transformation of $g$ into $\mathfrak{g}$; finally, from this, by transposition, derive the transformation of $\mathfrak{G}$ into $G$.}
shows that, if we set
\[
\begin{aligned}
(\beta'\gamma''-\beta''\gamma')\delta+
(\beta''\gamma-\beta\gamma'')\delta'+
(\beta\gamma'-\beta'\gamma)\delta''&=\zeta\\
(\gamma'\alpha''-\gamma''\alpha')\delta+
(\gamma''\alpha-\gamma\alpha'')\delta'+
(\gamma\alpha'-\gamma'\alpha)\delta''&=\eta
\end{aligned}
\]
the form $\mathfrak{G}$ is transformed into $G$ by the substitution
\[
\begin{array}{ccc}
k,&0,&0\\
0,&k,&0\\
\zeta,&\eta,&\mathfrak{k}
\end{array}
\]
This gives
\[
B=\eta\mathfrak{k}D+\mathfrak{k}k\mathfrak{B},\qquad
B'=\zeta\mathfrak{k}D+\mathfrak{k}k\mathfrak{B}'
\]
and so, since $\mathfrak{k}k=\pm1$, either $B\equiv\mathfrak{B}$, $B'\equiv\mathfrak{B}'$, or $B\equiv-\mathfrak{B}$, $B'\equiv-\mathfrak{B}'$
$(\mathrm{mod.}\ D)$. In the first case we call the values $(B,B')$, $(\mathfrak{B},\mathfrak{B}')$ equivalent;
in the second, opposite. We will say that the representation of the form
$\varphi$ \textit{belongs} to any value of the
\SourceLeaf{330}
expression $\surd\Delta(p,\allowbreak -q,\allowbreak r)(\mathrm{mod.}\ D)$ that can be derived from it by the method in I.
Thus all values to which the same representation belongs will be either
equivalent or opposite.

III.\ Conversely, suppose that, as in I above, the representation of the
form $\varphi$ by $f$ given by $x=\alpha t+\beta u$ and so on belongs to the value
$(B,B')$ derived from it with the help of the transformation
\[
\begin{array}{ccc}
\alpha,&\beta,&\gamma\\
\alpha',&\beta',&\gamma'\\
\alpha'',&\beta'',&\gamma''
\end{array}
\]
It will also belong to any other value $(\mathfrak{B},\mathfrak{B}')$ that is equivalent or
opposite to the first; \textit{i.e.}, in place of
$\gamma,\gamma',\gamma''$ we may choose other integers $\delta,\delta',\delta''$ for which the equation
$(\Omega)$, namely
\[
(\alpha'\beta''-\alpha''\beta')\delta+
(\alpha''\beta-\alpha\beta'')\delta'+
(\alpha\beta'-\alpha'\beta)\delta''=\pm1
\]
holds, and which are such that the coefficients numbered $4$ and
$5$ in the form adjoint to the one into which $f$ is transformed
by the substitution $(S)$
\[
\begin{array}{ccc}
\alpha,&\beta,&\delta\\
\alpha',&\beta',&\delta'\\
\alpha'',&\beta'',&\delta''
\end{array}
\]
become $=\mathfrak{B},\mathfrak{B}'$, respectively. For set
\[
\pm B=\mathfrak{B}+\eta D,\qquad \pm B'=\mathfrak{B}'+\zeta D
\]
(taking the upper or lower signs here and below according as the values
$(B.\ B')$, $(\mathfrak{B},\mathfrak{B}')$ are equivalent or opposite).
Then $\zeta,\eta$ will be integers. Let $g$ be transformed by the substitution
\[
\begin{array}{ccc}
1,&0,&\zeta\\
0,&1,&\eta\\
0,&0,&\pm1
\end{array}
\]
into the form $\mathfrak{g}$. It is easy to see that its determinant is
$\Delta$, and that in its adjoint form the coefficients numbered
$4$ and $5$ become $=\mathfrak{B},\mathfrak{B}'$, respectively. But on setting
\[
\alpha\zeta+\beta\eta\pm\gamma=\delta,\quad
\alpha'\zeta+\beta'\eta\pm\gamma'=\delta',\quad
\alpha''\zeta+\beta''\eta\pm\gamma''=\delta''
\]
\SourceLeaf{331}
it is easy to see that $f$ is transformed by the substitution
$(S)$ into $\mathfrak{g}$, and that the equation $(\Omega)$ is satisfied.
\textit{Q. E. D.}

